\id{IRSTI 82.17.252 }{}Sofia University `St. Kliment Ohridski'

{\bfseries EXPLAINING THE GAP BETWEEN INTEREST AND APPLICATION INTENTION IN
PUBLIC SECTOR EMPLOYMENT}

{\bfseries \tsp{1}M.Yerbatyrov}\fig{e/image1}[]{\bfseries ,\tsp{1}M.
Zhetpisbayeva}\fig{e/image1}[]{\bfseries ,
\tsp{2}Zh.
Vladimirov}\fig{e/image1}[]{\bfseries ,\tsp{3}G.
Baidauletova}\fig{e/image1}[]\envelope 

\emph{\tsp{1}Karaganda university of Kazpotrebsoyuz,
Karaganda, Kazakhstan,}

\emph{\tsp{2}} \emph{Sofia University «St. Kliment
Ohridski», Sofia, Bulgaria,}

\emph{\tsp{3}Shymkent University, Shymkent, Kazakhstan}

\envelope Corresponding author:
gulnur.baidauletova@mail.ru

Research on public sector attractiveness consistently reveals a gap
between positive evaluations of public missions and the actual
willingness of young people to apply for government positions. This
article aims to explain this gap by analytically distinguishing two
stages of career decision making: the formation of interest and the
formation of application intention. The empirical basis of the study
consists of survey data collected from students of Karaganda university
of KAZPOTREBSOYUZ (N = 196) in September 2025, who evaluated a real
vacancy in the public sector. The findings demonstrate that symbolic
attractiveness associated with mission and perceived public value plays
a significant role in generating interest, but loses explanatory power
at the stage of intention formation. Application intention emerges under
conditions of anticipated work constraints and early sectoral screening;
whereby innovation-oriented students exclude the public sector from
their set of considered alternatives. In addition, evaluative
ambivalence is identified as a stable condition that disrupts linear
transitions from values to action. The results refine the explanatory
scope of public service motivation (PSM) theory and highlight the need
to complement attraction models with mechanisms of early exclusion and
ambivalent evaluation.

{\bfseries Keywords:} public service motivation, public sector
attractiveness, application intention, innovation orientation,
ambivalence, early sectoral exclusion.

{\bfseries ОБЪЯСНЕНИЕ РАЗРЫВА МЕЖДУ ИНТЕРЕСОМ И НАМЕРЕНИЕМ ПОДАТЬ ЗАЯВКУ НА
РАБОТУ В ГОСУДАРСТВЕННОМ СЕКТОРЕ}

{\bfseries \tsp{1}М. Ербатыров, \tsp{1}М.
Жетписбаева, \tsp{2}Ж. Владимиров, \tsp{3}Г.О.
Байдаулетова\envelope }

\emph{\tsp{1}Карагандинский университет «Казпотребсоюза»,
Караганда, Казахстан,}

\emph{\tsp{2}Софийский университет «Св. Климент Охридского»,
София, Болгария,}

\emph{\tsp{3}Шымкентский университет, Шымкент, Казахстан,}

e-mail:
gulnur.baidauletova@mail.ru

В исследованиях привлекательности государственной службы устойчиво
фиксируется разрыв между положительной оценкой публичных миссий и
реальной готовностью молодых специалистов подавать заявления на
государственные должности. Цель настоящей статьи заключается в
объяснении данного разрыва через разграничение двух стадий карьерного
решения формирования интереса и формирования намерения действовать.
Эмпирическую основу исследования составляют данные анкетного опроса
студентов Карагандинского университета «КАЗПОТРЕБСОЮЗА» (N = 196),
проведённого в сентябре 2025 года при оценке конкретной вакансии в
государственном секторе. Показано, что символическая привлекательность,
связанная с миссией и общественной значимостью государственной службы,
оказывает существенное влияние на формирование интереса, однако
утрачивает объяснительную силу при переходе к намерению подать
заявление. Намерение формируется в условиях ожидаемых ограничений
повседневной работы и предварительного отраслевого отбора, в рамках
которого инновационно ориентированные студенты исключают государственный
сектор из числа рассматриваемых альтернатив. Кроме того, установлено,
что амбивалентность оценок представляет собой устойчивое состояние,
нарушающее линейный переход от ценностных установок к действию.
Полученные результаты уточняют границы применимости теории мотивации
государственного служения (PSM) и указывают на необходимость дополнения
моделей привлекательности механизмами раннего исключения и анализа
противоречивых оценок.

{\bfseries Ключевые слова:} мотивация государственной службы,
привлекательность государственной службы, намерение подать заявление,
инновационная ориентация, амбивалентность, раннее исключение сектора.

{\bfseries МЕМЛЕКЕТТІК СЕКТОРДАҒЫ ЖҰМЫСҚА ҚЫЗЫҒУШЫЛЫҚ ПЕН ӨТІНІШ БЕРУ
НИЕТІНІҢ АРАСЫНДАҒЫ АЛШАҚТЫҚТЫ ТҮСІНДІРУ}

{\bfseries \tsp{1}М. Ербатыров, \tsp{1}М.
Жетпісбаева, \tsp{2}Ж. Владимиров, \tsp{3}Г.О.
Байдаулетова\envelope }

\emph{\tsp{1}«Қазтұтынуодағы» Қарағанды университеті,
Қарағанды, Қазақстан,}

\emph{\tsp{2}Әулие Климент Охрид атындағы София
университеті, София, Болгария,}

\emph{\tsp{3}Шымкент университеті, Шымкент, Қазақстан}

e-mail:
gulnur.baidauletova@mail.ru

Мемлекеттік қызметтің тартымдылығын зерттейтін еңбектерде қоғамдық
миссияларды оң бағалау мен нақты мемлекеттік лауазымдарға өтініш беру
ниетінің арасындағы алшақтық тұрақты түрде байқалады. Осы мақаланың
мақсаты аталған алшақтықты мансаптық шешім қабылдаудың екі кезеңін, яғни
қызығушылықтың қалыптасуы мен нақты әрекетке ниеттің қалыптасуын ажырату
арқылы түсіндіру. Зерттеудің эмпирикалық негізін 2025 жылғы қыркүйек
айында «ҚАЗТҰТЫНУОДАҒЫ» Қарағанды университетінің студенттері арасында
(N = 196) жүргізілген сауалнама деректері құрады, онда студенттер
мемлекеттік сектордағы нақты бос жұмыс орнын бағалады. Зерттеу
нәтижелері мемлекеттік қызметтің миссиясы мен қоғамдық маңыздылығына
байланысты символдық тартымдылықтың қызығушылықтың қалыптасуына елеулі
әсер ететінін, алайда өтініш беру ниетіне көшу кезеңінде оның
түсіндірмелік әлеуеті әлсірейтінін көрсетеді. Ниет күнделікті еңбек
қызметіне қатысты күтілетін шектеулер мен алдын ала салалық іріктеу
жағдайында қалыптасады, бұл ретте инновацияға бағдарланған студенттер
мемлекеттік секторды қарастырылатын нұсқалар қатарынан ерте кезеңде алып
тастайды. Сонымен қатар бағалаулардың амбиваленттілігі құндылықтардан
нақты әрекетке өтудің сызықтық логикасын бұзатын тұрақты күй ретінде
айқындалды. Алынған нәтижелер мемлекеттік қызмет мотивациясы (PSM)
теориясының қолданылу шектерін нақтылап, тартымдылық модельдерін ерте
алып тастау механизмдерімен толықтыру қажеттігін көрсетеді.

{\bfseries Түйін сөздер:} мемлекеттік қызмет мотивациясы, мемлекеттік
қызметтің тартымдылығы, өтініш беру ниеті, инновацияға бағдарлану,
амбиваленттілік, саланы ерте алып тастау.

{\bfseries Introduction.} Public sector attraction research has become
unusually consistent in documenting one empirical pattern: early career
candidates often endorse public missions, yet frequently hesitate to
apply for public sector jobs. This pattern has been observed across
countries, cohorts, and institutional settings, and it persists even
where public sector work is evaluated positively in terms of social
contribution and legitimacy {[}1 - 3{]}. Rather than disappearing with
improved measurement or changing labor market conditions, the gap
between positive evaluation and application intention has become one of
the field's most stable findings {[}4{]}.

Dominant explanations of this gap are rooted in public service
motivation (PSM) and related attraction frameworks. These approaches
assume that symbolic alignment with public values such as commitment to
societal contribution, collective goals, or institutional purpose enters
the decision process early and exerts a downstream influence on
behavioral intention {[}5{]}. In our setting, weak application rates are
typically interpreted as the result of incomplete information, competing
incentives, or external constraints that interfere with otherwise
positive motivations. Implicitly, intention is modeled as the behavioral
endpoint of symbolic appeal {[}6{]}.

However, accumulating empirical evidence challenges this assumption.
Multiple studies report that once expectations about day-to-day work
realities autonomy, flexibility, pace of decision making, and
opportunities for initiative are taken into account, symbolic appeal no
longer predicts application behavior in a reliable or consistent way
{[}5{]}. Importantly, this pattern does not suggest that symbolic
motivations are irrelevant. Instead, it indicates that their influence
may be confined to earlier stages of evaluation and may not extend to
the point at which individuals decide whether to apply.

This paper argues that the persistence of the attraction intention gap
is not merely an empirical anomaly but a conceptual problem {[}7{]}.
Specifically, prevailing attraction models' mis specify the decision
process by treating application intention as a continuation of symbolic
interest. Intention is not formed in the same evaluative environment as
interest {[}8{]}. While interest reflects a broad assessment of meaning
and desirability, intention emerges under anticipated constraints,
sector level screening, and uncertainty about everyday work experiences.
Under these conditions, symbolic evaluations do not operate as direct
behavioral drivers {[}8{]}.

Two overlooked mechanisms are central to this mi's specification. The
first is early exclusion based on innovation related expectations. For
many students, preferences for flexibility, autonomy, and dynamic work
environments function as a preliminary filter that determines whether
the public sector enters the considered set at all. When such
expectations are salient, symbolic appeal does not compete with other
job attributes; it becomes behaviorally irrelevant because the sector
has already been screened out. In this study, innovation orientation
operates not as a tradeoff variable but as a sector level exclusion
logic.

The second mechanism is evaluative ambivalence. In practice than holding
coherent and internally consistent views, students frequently evaluate
public sector work in mixed terms simultaneously appreciating its
societal value while doubting its capacity to offer engaging or
autonomous work. This ambivalence is not a transient state or a
measurement artifact. It represents a stable evaluative condition under
which linear value to action models lose predictive power. When positive
and negative judgments coexist, symbolic appeal may generate interest
without resolving hesitation, leaving intention underdetermined.

Building on these arguments, this study reconceptualizes early career
attraction to public sector employment by analytically separating
interest from intention. Using survey data from students evaluating a
real public sector vacancy, the analysis examines how symbolic appeal,
innovation-oriented expectations, and ambivalence jointly shape
different stages of the decision process. The findings demonstrate that
while PSM and symbolic evaluations remain powerful predictors of
interest, they are conceptually insufficient for explaining why
individuals decide to apply. By locating the breakdown between
evaluation and action, the paper refocuses public sector attraction
theory on the conditions under which valued public service becomes
actionable {[}9{]}.

{\bfseries Materials and Methods.} The study uses a cross sectional survey
design to examine how symbolic appeal, innovation-oriented expectations,
and evaluative ambivalence shape distinct stages of early career
decision making: interest and application intention. Instead, then
treating attraction as a single latent construct, the design explicitly
separates evaluative appreciation from behavioral readiness, consistent
with the paper's theoretical argument.

To reduce hypothetical bias, respondents evaluated a real public sector
vacancy (Deputy head of a state institution, apparatus of the Akimat of
the Karagandy region). Anchoring evaluations in a concrete job
description ensures that responses reflect anticipated work realities,
not abstract sectoral images. The context is analytically relevant
because public sector employment in this setting is institutionally
salient and widely associated with organizational constraints, allowing
early exclusion mechanisms and ambivalence to surface clearly.

Data were collected between 17 and 24 September 2025 at Karaganda
university of «KAZPOTREBSOYUZ». A targeted non probability sampling
strategy was employed {[}6{]}. While this approach does not permit
population level prevalence estimates, it is appropriate for theory
driven research focused on identifying decision mechanisms instead than
estimating population parameters {[}7{]}.

The sample is predominantly female (74\%), with most respondents in
their third or fourth year of study. A majority reported no formal work
experience, while approximately two thirds had completed an internship
in a public organization. Gender and internship experience are included
as controls in all models due to their documented associations with job
interest and application behavior. An a priori power analysis indicated
that the final sample provides sufficient power to detect medium sized
effects in both linear and logistic regression models.

All variables were measured using 7-point Likert scales.

Interest and intention were measured with items capturing perceived job
attractiveness, willingness to accept an offer, initiative to apply, and
readiness to relocate if required.

Symbolic appeal was measured using items capturing mission alignment,
perceived societal contribution, and meaningfulness of public service.

Innovation orientation captured preferences for flexibility, autonomy,
and a dynamic work pace.

Ambivalence was operationalized using a cross valence index combining
positive and negative evaluations of public sector work; reliability was
assessed using the Spearman-Brown coefficient.

The analysis proceeded in three steps aligned with the paper's
theoretical claims. First, descriptive analyses assessed overall
evaluation patterns and divergence between interest and intention.
Second, linear regression models examined predictors of interest,
testing whether symbolic appeal remains dominant once innovation
expectations and ambivalence are included. Third, logistic regression
models examined application intention, treating intention as a
behavioral decision in practice than a continuous preference.

Predictor variables were mean centered prior to analysis, and
interaction terms were computed as multiplicative products. Intention
was dichotomized for logistic regression (0 = unlikely or neutral; 1 =
likely).

Participation was voluntary, informed consent was obtained, and
responses were anonymous. All procedures complied with institutional
ethical standards.

{\bfseries Results and Discussion.} Table 1 summarizes the characteristics
of the sample (N = 196). The sample is predominantly female (74\%), with
most respondents in their third or fourth year of study (92\%). Slightly
more than half of the respondents (54\%) receive a state scholarship. A
substantial majority report no formal work experience (75\%), while 67\%
have completed an internship in a public institution.

{\bfseries Table 1 - Sample characteristics (N = 196)}

%% \begin{longtable}[]{@{}
%%   >{\centering\arraybackslash}p{(\linewidth - 4\tabcolsep) * \real{0.3523}}
%%   >{\centering\arraybackslash}p{(\linewidth - 4\tabcolsep) * \real{0.4116}}
%%   >{\centering\arraybackslash}p{(\linewidth - 4\tabcolsep) * \real{0.2352}}@{}}
%% \toprule\noalign{}
%% \begin{minipage}[b]{\linewidth}\centering
%% {\bfseries Variable}
%% \end{minipage} & \begin{minipage}[b]{\linewidth}\centering
%% {\bfseries Categories}
%% \end{minipage} & \begin{minipage}[b]{\linewidth}\centering
%% {\bfseries \%}
%% \end{minipage} \\
%% \midrule\noalign{}
%% \endhead
%% \bottomrule\noalign{}
%% \endlastfoot
%% Gender & Female (145), Male (51) & 74 / 26 \\
%% Year of study & 1st (7), 2nd (8), 3rd (98), 4th (83) & 4 / 4 / 50 /
%% 42 \\
%% State scholarship & Yes (106), No (90) & 54 / 46 \\
%% Work experience & Yes (49), No (147) & 25 / 75 \\
%% Internship in public institutions & Yes (131), No (65) & 67 / 33 \\
%% \multicolumn{3}{@{}>{\raggedleft\arraybackslash}p{(\linewidth - 4\tabcolsep) * \real{0.9992} + 4\tabcolsep}@{}}{%
%% Note: Authors' calculations based on survey data} \\
%% \end{longtable}

This composition reflects a typical early career student population
evaluating first full-time employment opportunities. Gender and
internship experience are therefore included as control variables in all
subsequent model specifications.

Table 2 reports descriptive statistics for the three indicators of job
interest. Overall evaluations are moderately positive: mean values range
from 4.3 to 5.0 on the 7-point scale. Organizational attractiveness (M =
4,9) and willingness to accept an offer (M = 5,0) show relatively high
central tendency, with approximately 60\% of respondents selecting
values of 5 or higher.

{\bfseries Table 2 - Descriptive statistics for interest indicators}

%% \begin{longtable}[]{@{}
%%   >{\centering\arraybackslash}p{(\linewidth - 8\tabcolsep) * \real{0.3065}}
%%   >{\centering\arraybackslash}p{(\linewidth - 8\tabcolsep) * \real{0.1745}}
%%   >{\centering\arraybackslash}p{(\linewidth - 8\tabcolsep) * \real{0.1720}}
%%   >{\centering\arraybackslash}p{(\linewidth - 8\tabcolsep) * \real{0.1737}}
%%   >{\centering\arraybackslash}p{(\linewidth - 8\tabcolsep) * \real{0.1734}}@{}}
%% \toprule\noalign{}
%% \begin{minipage}[b]{\linewidth}\centering
%% {\bfseries Indicator}
%% \end{minipage} & \begin{minipage}[b]{\linewidth}\centering
%% {\bfseries Mean}
%% \end{minipage} & \begin{minipage}[b]{\linewidth}\centering
%% {\bfseries \% ≥5}
%% \end{minipage} & \begin{minipage}[b]{\linewidth}\centering
%% {\bfseries \% Max}
%% \end{minipage} & \begin{minipage}[b]{\linewidth}\centering
%% {\bfseries \% Min}
%% \end{minipage} \\
%% \midrule\noalign{}
%% \endhead
%% \bottomrule\noalign{}
%% \endlastfoot
%% Attractiveness & 4,9 & 62 & 29 & 9 \\
%% Accept offer & 5,0 & 60 & 42 & 9 \\
%% Apply independently & 4,3 & 49 & 28 & 25 \\
%% \multicolumn{5}{@{}>{\raggedleft\arraybackslash}p{(\linewidth - 8\tabcolsep) * \real{1.0000} + 8\tabcolsep}@{}}{%
%% Note: Authors' calculations based on survey data} \\
%% \end{longtable}

In contrast, initiative to apply independently exhibits a markedly
different distribution. Although its mean remains above the scale
midpoint (M = 4.3), responses are strongly polarized, with 28\%
selecting the maximum category and 25\% selecting the minimum category.

This divergence indicates that positive evaluations of the vacancy do
not translate uniformly into readiness to act, providing early
descriptive evidence for a separation between general interest and
application intention.

Gender differences are pronounced. Among respondents fully willing to
apply (score = 7), 93\% are women (51 of 55), whereas among respondents
fully unwilling to apply (score = 1), 71\% are men (35 of 49). Other
background characteristics show no systematic association with interest
indicators. Year of study, scholarship status, and prior work experience
do not meaningfully differentiate evaluations. Internship experience in
public institutions is descriptively associated with higher interest but
does not exhibit sufficient variation to support robust statistical
inference and is therefore treated as a control variable in subsequent
models.

Table 3 presents Spearman correlations between value orientations and
indicators of job interest. Two consistent patterns emerge. Orientations
toward social mission, stability, and structured work are positively
associated with all interest indicators, with correlations ranging from
ρ ≈ 0,32 to 0,52. In contrast, innovation orientation is negatively
associated with each indicator, with correlations ranging from ρ ≈
--0,27 to -0,55.

{\bfseries Table 3 - Spearman correlations (ρ) between values and interest
indicators}

%% \begin{longtable}[]{@{}
%%   >{\centering\arraybackslash}p{(\linewidth - 10\tabcolsep) * \real{0.2347}}
%%   >{\centering\arraybackslash}p{(\linewidth - 10\tabcolsep) * \real{0.1470}}
%%   >{\centering\arraybackslash}p{(\linewidth - 10\tabcolsep) * \real{0.1470}}
%%   >{\centering\arraybackslash}p{(\linewidth - 10\tabcolsep) * \real{0.1470}}
%%   >{\centering\arraybackslash}p{(\linewidth - 10\tabcolsep) * \real{0.1470}}
%%   >{\centering\arraybackslash}p{(\linewidth - 10\tabcolsep) * \real{0.1775}}@{}}
%% \toprule\noalign{}
%% \begin{minipage}[b]{\linewidth}\centering
%% {\bfseries Orientation}
%% \end{minipage} & \begin{minipage}[b]{\linewidth}\centering
%% {\bfseries Attr.}
%% \end{minipage} & \begin{minipage}[b]{\linewidth}\centering
%% {\bfseries Accept}
%% \end{minipage} & \begin{minipage}[b]{\linewidth}\centering
%% {\bfseries Apply}
%% \end{minipage} & \begin{minipage}[b]{\linewidth}\centering
%% {\bfseries Relocate}
%% \end{minipage} & \begin{minipage}[b]{\linewidth}\centering
%% {\bfseries Consider}
%% \end{minipage} \\
%% \midrule\noalign{}
%% \endhead
%% \bottomrule\noalign{}
%% \endlastfoot
%% Mission & 0,46 & 0,52 & 0,38 & 0,38 & 0,50 \\
%% Stability & 0,48 & 0,38 & 0,38 & 0,31 & 0,51 \\
%% Structure & 0,42 & 0,32 & 0,37 & 0,13 & 0,36 \\
%% Innovation & --0,55 & --0,47 & --0,41 & --0,27 & --0,53 \\
%% \multicolumn{6}{@{}>{\raggedleft\arraybackslash}p{(\linewidth - 10\tabcolsep) * \real{1.0000} + 10\tabcolsep}@{}}{%
%% Note: Authors' calculations based on survey data} \\
%% \end{longtable}

Notably, innovation orientation exhibits the strongest associations in
absolute terms, particularly with organizational attractiveness and
willingness to accept an offer. This pattern indicates that preferences
for autonomy and dynamic work environments are systematically related to
lower evaluations of the public sector vacancy at the interest stage.

Overall, the correlation structure shows a clear alignment between value
orientations and early evaluative judgments, while remaining agnostic
about whether these associations extend to application intention a
question addressed in subsequent multivariate analyses.

Table 4 reports the results of an exploratory factor analysis conducted
on the 18 attitude items. Sampling adequacy was high (KMO = 0,82),
exceeding commonly accepted thresholds for factorability. Three factors
with eigenvalues greater than one were extracted, jointly explaining
68\% of the total variance.

{\bfseries Table 4 - Exploratory factor analysis of attitudinal items (N =
196)}

%% \begin{longtable}[]{@{}
%%   >{\centering\arraybackslash}p{(\linewidth - 6\tabcolsep) * \real{0.2809}}
%%   >{\centering\arraybackslash}p{(\linewidth - 6\tabcolsep) * \real{0.1949}}
%%   >{\centering\arraybackslash}p{(\linewidth - 6\tabcolsep) * \real{0.1823}}
%%   >{\centering\arraybackslash}p{(\linewidth - 6\tabcolsep) * \real{0.3419}}@{}}
%% \toprule\noalign{}
%% \begin{minipage}[b]{\linewidth}\centering
%% {\bfseries Item}
%% \end{minipage} & \begin{minipage}[b]{\linewidth}\centering
%% {\bfseries General Evaluation}
%% \end{minipage} & \begin{minipage}[b]{\linewidth}\centering
%% {\bfseries Mission \& Prestige}
%% \end{minipage} & \begin{minipage}[b]{\linewidth}\centering
%% {\bfseries Material \& Organizational Conditions}
%% \end{minipage} \\
%% \midrule\noalign{}
%% \endhead
%% \bottomrule\noalign{}
%% \endlastfoot
%% Organizational attractiveness & 0,81 & - & - \\
%% Willingness to accept an offer & 0,78 & - & - \\
%% Initiative to apply & 0,72 & - & - \\
%% Recommend the job to others & 0,69 & - & - \\
%% Societal contribution & - & 0,82 & - \\
%% Public value & - & 0,79 & - \\
%% Trustworthiness of organization & - & 0,74 & - \\
%% Organizational prestige & - & 0,71 & - \\
%% Salary attractiveness & - & - & 0,76 \\
%% Job stability & - & - & 0,73 \\
%% Working conditions & - & - & 0,68 \\
%% \multicolumn{4}{@{}>{\raggedleft\arraybackslash}p{(\linewidth - 6\tabcolsep) * \real{1.0000} + 6\tabcolsep}@{}}{%
%% Note: Authors' calculations based on survey data} \\
%% \end{longtable}

The first factor, general evaluation, captures overall positive
assessment of the vacancy, including items related to organizational
attractiveness, willingness to accept an offer, willingness to apply,
and recommending the job. The second factor, mission and prestige,
reflects symbolic evaluations such as societal contribution, public
value, trustworthiness, and perceived prestige. The third factor,
material and organizational conditions, captures evaluations of salary,
stability, and working conditions. Factor loadings ranged from 0,61 to
0,82 across all factors.

Although application related items load together with general evaluative
items at the measurement level, subsequent analyses treat interest and
application intention as analytically distinct outcomes, consistent with
the theoretical argument that intention is formed under different
decision conditions than general evaluation.

Factor scores for mission and prestige and for material and
organizational conditions were used as predictors in subsequent
regression analyses.

Table 5 presents results from a multiple linear regression predicting
overall interest in the public sector vacancy. The dependent variable is
an interest index, computed as the mean of organizational
attractiveness, willingness to accept an offer, and initiative to apply.

{\bfseries Table 5- Multiple linear regression predicting overall interest}

%% \begin{longtable}[]{@{}
%%   >{\centering\arraybackslash}p{(\linewidth - 6\tabcolsep) * \real{0.4115}}
%%   >{\centering\arraybackslash}p{(\linewidth - 6\tabcolsep) * \real{0.1766}}
%%   >{\centering\arraybackslash}p{(\linewidth - 6\tabcolsep) * \real{0.1766}}
%%   >{\centering\arraybackslash}p{(\linewidth - 6\tabcolsep) * \real{0.2354}}@{}}
%% \toprule\noalign{}
%% \begin{minipage}[b]{\linewidth}\centering
%% {\bfseries Predictor}
%% \end{minipage} & \begin{minipage}[b]{\linewidth}\centering
%% {\bfseries Std. β}
%% \end{minipage} & \begin{minipage}[b]{\linewidth}\centering
%% {\bfseries SE}
%% \end{minipage} & \begin{minipage}[b]{\linewidth}\centering
%% {\bfseries P value}
%% \end{minipage} \\
%% \midrule\noalign{}
%% \endhead
%% \bottomrule\noalign{}
%% \endlastfoot
%% Mission \& prestige & 0,43 & 0,06 & <{} 0,001 \\
%% Material \& organizational conditions & 0,25 & 0,07 & <{}
%% 0,01 \\
%% Gender (female = 1) & 0,19 & 0,08 & <{} 0,05 \\
%% Internship experience & 0,07 & 0,06 & n.s. \\
%% \multicolumn{4}{@{}>{\raggedleft\arraybackslash}p{(\linewidth - 6\tabcolsep) * \real{1.0000} + 6\tabcolsep}@{}}{%
%% Note: Authors' calculations based on survey data} \\
%% \end{longtable}

Mission and prestige emerge as the strongest predictor of interest (β ≈
0,43, p <{} 0,001). Material and organizational conditions also
contribute positively, though more moderately (β ≈ 0,25, p <{}
0,01). Gender remains significant net of evaluative controls (β ≈ 0,19,
p <{} 0,05), while internship experience does not add
explanatory power once subjective evaluations are included. The model
explains approximately 45\% of the variance in interest (R² ≈ 0,45).

These results indicate that symbolic evaluations dominate the formation
of interest, whereas background characteristics play a secondary role
once perceptions of the job are taken into account.

Table 6 reports results from a logistic regression model estimating the
probability of intending to apply for the vacancy. The dependent
variable is binary and coded as 1 if the respondent reports an intention
to apply (score ≥ 5) and 0 otherwise, consistent with prior research
treating application intention as a threshold decision in practice than
a continuous preference.

{\bfseries Table 6 - Logistic regression estimates of application intention
(N = 196)}

%% \begin{longtable}[]{@{}
%%   >{\centering\arraybackslash}p{(\linewidth - 6\tabcolsep) * \real{0.3822}}
%%   >{\centering\arraybackslash}p{(\linewidth - 6\tabcolsep) * \real{0.2355}}
%%   >{\centering\arraybackslash}p{(\linewidth - 6\tabcolsep) * \real{0.1766}}
%%   >{\centering\arraybackslash}p{(\linewidth - 6\tabcolsep) * \real{0.2056}}@{}}
%% \toprule\noalign{}
%% \begin{minipage}[b]{\linewidth}\centering
%% {\bfseries Predictor}
%% \end{minipage} & \begin{minipage}[b]{\linewidth}\centering
%% {\bfseries Odds Ratio}
%% \end{minipage} & \begin{minipage}[b]{\linewidth}\centering
%% {\bfseries SE}
%% \end{minipage} & \begin{minipage}[b]{\linewidth}\centering
%% {\bfseries P value}
%% \end{minipage} \\
%% \midrule\noalign{}
%% \endhead
%% \bottomrule\noalign{}
%% \endlastfoot
%% Mission \& Prestige & 1.45 & 0.18 & <{} .01 \\
%% Stability orientation & 1.32 & 0.16 & <{} .05 \\
%% Innovation orientation & 0.68 & 0.14 & <{} .01 \\
%% Gender (female = 1) & 1.60 & 0.25 & <{} .01 \\
%% Internship experience & 1.18 & 0.20 & n.s. \\
%% \multicolumn{4}{@{}>{\raggedleft\arraybackslash}p{(\linewidth - 6\tabcolsep) * \real{1.0000} + 6\tabcolsep}@{}}{%
%% Note: Authors' calculations based on survey data} \\
%% \end{longtable}

The results reveal a clear and asymmetric pattern. Higher evaluations of
mission and Prestige are associated with an increased probability of
applying, indicating that symbolic appeal continues to matter at the
threshold of action. Stability orientation also contributes positively
to application likelihood. In contrast, innovation orientation is
associated with a lower probability of applying, consistent with an
early exclusion mechanism. Gender differences persist net of evaluative
controls, with women showing a significantly higher probability of
intending to apply. Internship experience in public institutions has a
positive but comparatively weaker association with application
intention.

Overall, the logistic regression results indicate that while symbolic
and stability related evaluations increase the likelihood of applying,
innovation-oriented expectations systematically reduce it. This pattern
mirrors the descriptive distributions and reinforces the distinction
between general interest and application intention.

This study set out to explain a persistent paradox in public sector
attraction research: why positive evaluations of public missions so
often fail to translate into application intentions. The results provide
a clear answer. The gap between interest and intention is not a matter
of weak motivation, missing incentives, or insufficient information. It
reflects a misalignment between how attraction has been theorized and
how early career decisions are actually formed.

The findings confirm a robust and well documented pattern: symbolic
appeal mission, societal contribution, and perceived public value
strongly predicts interest in public sector work. Students who value
public service express higher attractiveness ratings and greater
openness toward public sector employment. In this study, the empirical
foundations of PSM theory remain intact.

At the intention stage, however, symbolic appeal does not carry over to
the intention stage once anticipated work constraints are taken into
account. When expectations about autonomy, flexibility, and everyday
work dynamics enter the decision process, symbolic evaluations lose
their predictive relevance for application behavior. Symbolic
motivations continue to matter, though their effect appears limited to
early evaluative stages. In practice, it demonstrates that their
influence is structurally confined to the formation of interest.
Treating intention as a behavioral extension of symbolic alignment
therefore overstates the scope of PSM and obscures where the decision
process actually breaks down.

A central contribution of this study lies in identifying innovation
orientation as an early exclusion mechanism. Students who associate
themselves with dynamic, flexible, and innovation driven work
environments are significantly less likely to consider public sector
employment, not because they reject its mission, but because they screen
out the sector before symbolic considerations can become behaviorally
relevant.

For many respondents, especially those with strong innovation
preferences, this pattern challenges attraction models based on
attribute tradeoffs. In the present case, innovation orientation does
not compete with symbolic appeal; it precedes it. Once the public sector
is excluded from the considered set, symbolic appeal no longer functions
as a motivational input. This exclusion logic helps explain why public
sector recruitment efforts that emphasize mission and societal impact
often fail to reach innovation-oriented candidates, even when these
candidates express strong prosocial values.

The results also demonstrate that ambivalence is a defining feature of
how students evaluate public sector work. Instead, then expressing
coherent or one-sided attitudes, many respondents simultaneously
acknowledge the societal value of public service and doubt its capacity
to offer engaging or autonomous work experiences. This ambivalence is
not a transient hesitation or a measurement artifact. It represents a
stable evaluative state in which positive and negative judgments
coexist.

Under conditions of ambivalence, linear value to action models lose
explanatory power. Symbolic appeal can increase interest without
resolving hesitation, leaving intention underdetermined. From this
perspective, the attraction intention gap is not surprising; it is a
predictable outcome of decision making under evaluative conflict. Models
that assume evaluative coherence therefore misinterpret ambivalence as
noise, when it is in fact a core feature of the decision environment.

Taken together, these findings imply that public sector attraction
theory needs to be reformulated around a stage sensitive view of
decision making. Interest and intention are not successive points on a
single continuum but outcomes of distinct evaluative processes. Symbolic
appeal shapes interest by providing meaning and legitimacy. Innovation
orientation determines whether the public sector is considered at all.
Ambivalence characterizes the context in which intention is formed and
explains why interest frequently fails to convert into action.

From this perspective, the limitations of PSM theory are conceptual
rather than empirical {[}5{]}. PSM explains why individuals like public
sector work, but it cannot, on its own, explain why they decide to apply
{[}8{]}. Addressing this limitation requires moving beyond linear
attraction models and theorizing intention as a decision formed under
anticipated constraint and unresolved evaluative conflict.

For research, these findings call for a clearer analytical separation
between interest and intention in studies of public sector attraction.
Treating intention as a distinct decision outcome opens new avenues for
theorizing exclusion mechanisms, ambivalence, and sector level screening
processes that are obscured in conventional models.

For practice, the results suggest that recruitment strategies focused
exclusively on mission and symbolic appeal are unlikely to convert
interest into applications among early career candidates. Addressing
anticipated work constraints particularly perceptions of autonomy and
innovation may be necessary not to increase symbolic appeal, but to
prevent early exclusion from the considered set.

Overall, the study shows that the public sector attracts through meaning
but loses candidates through anticipated constraint. Understanding this
dynamic requires rethinking attraction not as a linear motivational
process, but as a sequence of evaluative filters that shape whether
valued public service becomes actionable.

{\bfseries Conclusion.} This paper addressed a persistent paradox in public
sector attraction research: students endorse public missions and report
high symbolic appreciation, yet often do not intend to apply for public
sector jobs. The findings show that this paradox is not a puzzle of weak
motivation or inadequate incentives. It is a predictable outcome of a
mis specified decision model one that treats application intention as
the behavioral continuation of symbolic appeal.

Three conclusions follow.

First, symbolic appeal is a reliable mechanism of interest formation,
not of intention. PSM and mission related evaluations continue to
explain why early career candidates like public sector work. Yet once
anticipated work constraints enter the decision especially expectations
about autonomy, flexibility, and the dynamics of daily work symbolic
evaluations no longer determine whether candidates intend to apply. The
theoretical implication is not that PSM «fails», but that its
explanatory scope is stage bound: it accounts for attraction as meaning,
not for intention as action.

Second, innovation orientation operates as a sector level exclusion
logic. For a substantial segment of students, preferences for dynamic
and flexible work environments function as an early screening criterion
that removes the public sector from the considered set before symbolic
considerations can become behaviorally relevant. This mechanism helps
explain why mission focused recruitment frequently produces interest
without conversion: the candidates most sensitive to autonomy and
innovation are filtered out prior to attribute tradeoffs.

Third, ambivalence is not noise but a stable evaluative state that
reshapes how intention is formed. Students often hold mixed judgments
valuing the societal contribution of public service while doubting its
capacity to deliver engaging and autonomous work. Under such conditions,
linear value to action models break down: symbolic appeal may raise
interest, but it does not resolve hesitation, leaving intention
underdetermined {[}9{]}.

Taken together, these conclusions require a reformulation of public
sector attraction theory. Interest and intention should not be treated
as successive points on a single continuum but as outcomes of distinct
evaluative processes {[}10{]}. Attraction begins with symbolic meaning,
but application intention is formed under anticipated constraint, early
sectoral screening, and persistent ambivalence. Models that continue to
treat intention as downstream symbolic attraction will systematically
overestimate the behavioral reach of PSM.

For public employers, the practical implication is straightforward:
mission-based messaging can generate interest, but it cannot by itself
produce applications when candidates anticipate constrained autonomy and
limited opportunities for initiative. Recruitment strategies that aim to
convert early career interest into action must address the work dynamics
expectations that trigger early exclusion and sustain ambivalence.

Ultimately, this study changes what the field should treat as the core
explanatory problem. The question is not only who values public service,
but when valuing public service becomes actionable and when, by design
of the decision process, it does not.

{\bfseries References}

1. Gupta R., Kakkar S., Chaudhuri M., Malik P., Saha S. Motivated to
serve, conflicted to commit? Reassessing career satisfaction and
publicness in the public service motivation--commitment nexus. //Asia
Pacific Journal of Public Administration. -2026.

\href{https://doi.org/10.1080/23276665.2026.2618128}{DOI
10.1080/23276665.2026.2618128}.

2. Perry J. L., Hondeghem A., Wise L. R. Revisiting the motivational
bases of public service: Twenty years of research and an agenda for the
future. //Administration review. -2010.-Vol.70(5). -P.681-690.
\href{https://doi.org/10.1111/j.1540-6210.2010.02196.x}{DOI
10.1111/j.1540-6210.2010.02196.x}.

3. Vandenabeele W., Breaugh J. It takes two to tango: concepts and
evidence of further integration of public service motivation theory and
self-determination theory//International Public Management Journal. -
2025. -Vol.28(1). - P.133-152.
\href{https://doi.org/10.1080/10967494.2024.2386292}{DOI
10.1080/10967494.2024.2386292}.

4. Doverspike D., Flores C., VanderLeest J. Charter 14- Lifespan
perspectives on personnel selection and recruitment. //Work across the
lifespan.-2019.-P.343-368 DOI
\href{https://doi.org/10.1016/B978-0-12-812756-8.00014-1}{10.1016/B978-0-12-812756-8.00014-1}.

5. Vandenabeele W., Ritz A., Neumann O. Public service motivation: State
of the art and conceptual cleanup. //The Palgrave handbook of public
administration and management in Europe. -2017. -P.261-278.
\href{https://doi.org/10.1057/978-1-137-55269-3_13}{DOI
10.1057/978-1-137-55269-3\_13}.

6. Belle N., Cantarelli, P. Monetary incentives, motivation, and job
effort in the public sector: An experimental study with Italian
government executives. //Review of Public Personnel Administration.
-2015. Vol.35(2). - P.99-123.
\href{https://doi.org/10.1177/0734371X13520460}{DOI
10.1177/0734371X13520460}.

7. Cash P. J. (2018). Developing theory-driven design research//Design
Studies. - 2018.-Vol.56.-P.84 -119.
\href{https://doi.org/10.1016/j.destud.2018.03.002}{DOI
10.1016/j.destud.2018.03.002}.

8. Schott C., Neumann O., Bärtschi M., Ritz A. Public service motivation,
prosocial motivation, prosocial behavior, and altruism: Towards
disentanglement and conceptual clarity. //International

jurnal of Public Administration. -2019.- P.1200 - 1211. DOI
10.1080.01900692.2019.1588302.9.Nina Vari van Loon, W. Vandenabeele, P. Leisink. Clarifying the
relationship between public service motivation and in-role and
extra-role behaviors: The relative contributions of person-job and
person-organization fit. //The American Review of Public Administration.
-2017. Vol.47(6). -P.699-713.
\href{https://doi.org/10.1177/0275074015617547}{DOI
10.1177/0275074015617547}\ul{.}

10. Eslami, Z., \& Taheri, Y. S. Exploring the relationship between
public service motivation, prosocial motivation, and civil servants'
organizational citizenship Behaviour: findings from a systematic
literature review. -URL.
https://www.rechtswissenschaft.unibe.ch/unibe/portal/fak\_rechtwis/content/e6024/e6025/e118744/e1190006/e1416330/pane1727015/e1416332/files1416356/MA\_Eslami\_Taheri\_23.01.2023\_ger.pdf.-

Date of access: 20.12.2025.

\emph{{\bfseries Information about the authors}}

Yerbatyrov M.- PhD student, Karaganda university of Kazpotrebsoyuz,
Karaganda, Kazakhstan, e-mail:
yerbatyr.meirambek@gmail.com;

Zhetpisbayeva M.- Candidate of Economic Sciences, Professor, Karaganda
university of Kazpotrebsoyuz, Karaganda, Kazakhstan, e-mail:
rimakeu@mail.ru;

Vladimirov Zh.- Doctor of Economic Sciences, Professor, University of
Sofia St Kliment Ohridski, Sofia, Bulgaria, e-mail:
jeve@feb-uni.sofia.bg;

Baidauletova G.- Master, senior lecturer, Shymkent University, Shymkent,
Kazakhstan, e-mail:
super\_khan@mail.ru.

\emph{{\bfseries Сведения об авторах}}

Ербатыров М.- PhD докторант, Карагандинский университет Казпотребсоюза,
Караганда, Казахстан, e-mail:
yerbatyr.meirambek@gmail.com;

Жетписбаева М.- кандидат экономических наук, профессор, Карагандинский
университет Казпотребсоюза, Караганда, Казахстан, e-mail:
rimakeu@mail.ru;

Владимиров Ж.- доктор экономических наук, профессор, Софийский
университет "Св. Климент Охридского", София, Болгария, e-mail:
jeve@feb-uni.sofia.bg;

Байдаулетова Г.- Магистр, старший преподаватель, Шымкентский
университет, Шымкент, Казахстан, e-mail:
super\_khan@mail.ru.\